%======================================================================
%  Capitulo 6 - Desarrollo proyecto de bases de datos (generado por gen/cap6.py)
%======================================================================

\section{Desarrollo proyecto de bases de datos}

Este capítulo define cómo se usará la base de datos del módulo de
operación de envíos: expresa en álgebra relacional y en SQL las
consultas que necesitan sus usuarios, y especifica las tablas,
relaciones, transacciones e informes. Del prototipo funcional se define
por ahora su arquitectura; su construcción, donde se crea la base de
datos y se ejecutan las consultas, se desarrollará en la siguiente
entrega.

\subsection{Algebra relacional de consultas MySQL}
\label{sec:algebra}

El álgebra relacional describe una consulta como una secuencia de
operaciones sobre relaciones, independiente del gestor de bases de
datos; SQL es la forma de escribirla para que el gestor la ejecute. En
esta sección se expresan en álgebra relacional y en SQL diez consultas
que responden a necesidades de los usuarios del módulo
(sección~\ref{sec:usuarios}). Entre todas usan las operaciones del
álgebra que resume la tabla~\ref{tab:operadores}.

La notación es la del texto guía del curso\footnote{A. Silberschatz,
H. F. Korth y S. Sudarshan, \textit{Fundamentos de bases de datos},
5.ª ed., McGraw-Hill, 2006, capítulos 2 y 3.}. Las relaciones se
escriben con el nombre de su tabla en mayúsculas, como en el esquema
relacional de la sección~\ref{sec:esquemarel}. La reunión natural
($\bowtie$) combina dos relaciones por sus atributos de igual nombre;
en este modelo esos atributos son las claves foráneas, que llevan el
mismo nombre que la clave primaria a la que apuntan. Cuando dos
relaciones comparten además otro atributo (por ejemplo, nombre o
estado), se renombra o se proyecta antes de reunirlas, y cuando la
clave foránea tiene otro nombre se usa una reunión theta. Para que las
expresiones se lean con facilidad, se escriben por pasos: cada paso
guarda un resultado intermedio con la asignación ($\leftarrow$) y el
último da la respuesta.

Cada consulta se escribe también en SQL estándar, válido tanto en
MySQL, que nombra la plantilla del proyecto, como en Oracle, el gestor
elegido en el capítulo 5. Los valores entre comillas (un número de
guía, un despacho, un rango de fechas) son ejemplos de los datos de
entrada que el usuario indica en cada caso. Su ejecución sobre la base
de datos hace parte del prototipo funcional
(sección~\ref{sec:prototipo}).

{\singlespacing\small
\begin{longtable}{L{2.4cm} L{2.4cm} L{5.6cm} L{4.0cm}}
\caption{Operaciones del álgebra relacional usadas en las consultas.}
\label{tab:operadores}\\
\toprule
\textbf{Notación} & \textbf{Operación} & \textbf{Resultado} & \textbf{Equivalente en SQL} \\
\midrule
\endfirsthead
\toprule
\textbf{Notación} & \textbf{Operación} & \textbf{Resultado} & \textbf{Equivalente en SQL} \\
\midrule
\endhead
\bottomrule
\endfoot
$\sigma_{P}(r)$ & Selección & Tuplas de r que cumplen el predicado P & WHERE \\
$\Pi_{A}(r)$ & Proyección & Atributos A de r, sin tuplas repetidas & SELECT A \\
$r \bowtie s$ & Reunión natural & Combina las tuplas de r y s que coinciden en los atributos de igual nombre & JOIN \dots\ ON (atributos comunes) \\
$r \bowtie_{\theta} s$ & Reunión theta & Combina las tuplas de r y s que cumplen la condición $\theta$ & JOIN \dots\ ON $\theta$ \\
$r \cup s$ & Unión & Tuplas que están en r o en s & UNION \\
$r - s$ & Diferencia & Tuplas de r que no están en s & NOT IN, NOT EXISTS, MINUS \\
$r \div s$ & División & Valores de r que se combinan con todas las tuplas de s & GROUP BY con HAVING COUNT \\
$_{G}\mathcal{G}_{F}(r)$ & Agregación & Agrupa r por los atributos G y calcula las funciones F (count, max, \dots) & GROUP BY, funciones de grupo \\
$\rho_{x(A)}(E)$ & Renombre & El resultado de E con nombre x y atributos A & Alias (AS) \\
$x \leftarrow E$ & Asignación & Guarda el resultado de E con el nombre x & Subconsulta o vista \\
\end{longtable}
\par}

\subsubsection{C1. Datos de un envío}
\label{q:C1}
\textit{Usuarios:} U3. \textit{Operaciones:} selección y proyección. La mesa de servicio consulta los datos de un envío a partir de su número de guía.

\begin{center}\adjustbox{max width=\linewidth}{$\displaystyle\begin{array}{l}E \leftarrow \sigma_{\text{numero\_guia} = \text{'EE0000000000003'}}(\text{ENVIO}) \\[4pt] \Pi_{\text{numero\_guia, fecha\_registro, fecha\_compromiso, peso\_real, fragil, alto\_valor}}(E)\end{array}$}\end{center}

\begin{lstlisting}[style=sqlpro]
SELECT numero_guia, fecha_registro, fecha_compromiso,
       peso_real, fragil, alto_valor
FROM   envio
WHERE  numero_guia = 'EE0000000000003'
\end{lstlisting}


\subsubsection{C2. Recorrido de un envío}
\label{q:C2}
\textit{Usuarios:} U3, U9. \textit{Operaciones:} selección, reunión natural y proyección. Muestra los eventos de un envío en orden, con el nombre del estado y el empleado responsable.

\begin{center}\adjustbox{max width=\linewidth}{$\displaystyle\begin{array}{l}E \leftarrow \sigma_{\text{numero\_guia} = \text{'EE0000000000003'}}(\text{EVENTO\_ENVIO}) \\[4pt] R \leftarrow E \bowtie \text{ESTADO\_ENVIO} \bowtie \text{EMPLEADO} \\[4pt] \Pi_{\text{fecha\_hora, nombre, nombres, observacion}}(R)\end{array}$}\end{center}

\begin{lstlisting}[style=sqlpro]
SELECT ev.fecha_hora, es.nombre AS estado,
       em.nombres AS responsable, ev.observacion
FROM   evento_envio ev
JOIN   estado_envio es ON es.id_estado   = ev.id_estado
JOIN   empleado em     ON em.id_empleado = ev.id_empleado
WHERE  ev.numero_guia = 'EE0000000000003'
ORDER  BY ev.fecha_hora
\end{lstlisting}
El orden (ORDER BY) no es una operación del álgebra relacional, porque una relación es un conjunto sin orden; se aplica al presentar el resultado.


\subsubsection{C3. Planilla de un despacho}
\label{q:C3}
\textit{Usuarios:} U2, U4. \textit{Operaciones:} renombre, selección, reunión natural y proyección. Lista los envíos asignados a un despacho en el orden de visita, con el destinatario, la dirección y la zona.

\begin{center}\adjustbox{max width=\linewidth}{$\displaystyle\begin{array}{l}Z \leftarrow \rho_{\text{Z(id\_zona, id\_centro, zona, localidades)}}(\text{ZONA}) \\[4pt] R \leftarrow \sigma_{\text{id\_despacho} = \text{15}}(\text{ASIGNACION\_DESPACHO}) \bowtie \text{ENVIO} \bowtie \text{DESTINATARIO} \bowtie Z \\[4pt] \Pi_{\text{orden\_visita, numero\_guia, nombre, direccion\_destino, zona, resultado}}(R)\end{array}$}\end{center}

\begin{lstlisting}[style=sqlpro]
SELECT a.orden_visita, a.numero_guia, d.nombre AS destinatario,
       e.direccion_destino, z.nombre AS zona, a.resultado
FROM   asignacion_despacho a
JOIN   envio e        ON e.numero_guia     = a.numero_guia
JOIN   destinatario d ON d.id_destinatario = e.id_destinatario
JOIN   zona z         ON z.id_zona         = e.id_zona
WHERE  a.id_despacho = 15
ORDER  BY a.orden_visita
\end{lstlisting}
ZONA y DESTINATARIO tienen ambos un atributo nombre; si no se renombra uno de ellos, la reunión natural exigiría que fueran iguales. En SQL ese papel lo cumplen los alias.


\subsubsection{C4. Envíos clasificados sin asignar}
\label{q:C4}
\textit{Usuarios:} U2. \textit{Operaciones:} selección, reunión natural, proyección y diferencia. Encuentra los envíos que ya fueron clasificados pero todavía no se han asignado a ningún despacho de entrega.

\begin{center}\adjustbox{max width=\linewidth}{$\displaystyle\begin{array}{l}C \leftarrow \Pi_{\text{numero\_guia}}(\text{EVENTO\_ENVIO} \bowtie \sigma_{\text{nombre} = \text{'CLASIFICADO'}}(\text{ESTADO\_ENVIO})) \\[4pt] S \leftarrow \Pi_{\text{numero\_guia}}(\text{ASIGNACION\_DESPACHO}) \\[4pt] C - S\end{array}$}\end{center}

\begin{lstlisting}[style=sqlpro]
SELECT ev.numero_guia
FROM   evento_envio ev
JOIN   estado_envio es ON es.id_estado = ev.id_estado
WHERE  es.nombre = 'CLASIFICADO'
AND    ev.numero_guia NOT IN (SELECT numero_guia
                              FROM   asignacion_despacho)
\end{lstlisting}
La diferencia se escribe con NOT IN, que funciona igual en MySQL y en Oracle; Oracle también la ofrece con el operador MINUS.


\subsubsection{C5. Envíos por zona y modalidad}
\label{q:C5}
\textit{Usuarios:} U1. \textit{Operaciones:} selección, renombre, reunión natural y agregación. Cuenta los envíos registrados en un periodo por zona de destino y modalidad.

\begin{center}\adjustbox{max width=\linewidth}{$\displaystyle\begin{array}{l}Z \leftarrow \rho_{\text{Z(id\_zona, id\_centro, zona, localidades)}}(\text{ZONA}) \\[4pt] M \leftarrow \rho_{\text{M(id\_modalidad, modalidad, plazo\_dias\_habiles)}}(\text{MODALIDAD}) \\[4pt] R \leftarrow \sigma_{\text{fecha\_registro} \geq \text{'2026-11-01'} \wedge \text{fecha\_registro} \leq \text{'2026-11-30'}}(\text{ENVIO}) \bowtie Z \bowtie M \\[4pt] {}_{\text{zona, modalidad}}\mathcal{G}_{\text{count}(\text{numero\_guia}) \text{ as } \text{envios}}(R)\end{array}$}\end{center}

\begin{lstlisting}[style=sqlpro]
SELECT z.nombre AS zona, m.nombre AS modalidad,
       COUNT(e.numero_guia) AS envios
FROM   envio e
JOIN   zona z      ON z.id_zona      = e.id_zona
JOIN   modalidad m ON m.id_modalidad = e.id_modalidad
WHERE  e.fecha_registro BETWEEN DATE '2026-11-01' AND DATE '2026-11-30'
GROUP  BY z.nombre, m.nombre
ORDER  BY z.nombre, m.nombre
\end{lstlisting}
ZONA y MODALIDAD también comparten el atributo nombre, por eso ambas se renombran antes de reunirlas.


\subsubsection{C6. Mensajeros con varias novedades}
\label{q:C6}
\textit{Usuarios:} U1. \textit{Operaciones:} selección, reunión natural, reunión theta y agregación. Identifica los mensajeros que registraron dos o más novedades en un periodo, como alerta para revisar sus rutas.

\begin{center}\adjustbox{max width=\linewidth}{$\displaystyle\begin{array}{l}N \leftarrow \Pi_{\text{id\_novedad, id\_despacho}}(\text{NOVEDAD} \bowtie \text{ASIGNACION\_DESPACHO}) \\[4pt] D \leftarrow \Pi_{\text{id\_novedad, id\_mensajero}}(N \bowtie \sigma_{\text{fecha} \geq \text{'2026-11-01'} \wedge \text{fecha} \leq \text{'2026-11-30'}}(\text{DESPACHO})) \\[4pt] R \leftarrow D \bowtie_{\text{id\_mensajero = id\_empleado}} \text{EMPLEADO} \\[4pt] T \leftarrow {}_{\text{nombres, apellidos}}\mathcal{G}_{\text{count}(\text{id\_novedad}) \text{ as } \text{novedades}}(R) \\[4pt] \sigma_{\text{novedades} \geq 2}(T)\end{array}$}\end{center}

\begin{lstlisting}[style=sqlpro]
SELECT em.nombres, em.apellidos,
       COUNT(n.id_novedad) AS novedades
FROM   novedad n
JOIN   asignacion_despacho a ON a.id_asignacion = n.id_asignacion
JOIN   despacho d            ON d.id_despacho   = a.id_despacho
JOIN   empleado em           ON em.id_empleado  = d.id_mensajero
WHERE  d.fecha BETWEEN DATE '2026-11-01' AND DATE '2026-11-30'
GROUP  BY em.nombres, em.apellidos
HAVING COUNT(n.id_novedad) >= 2
\end{lstlisting}
NOVEDAD y DESPACHO tienen ambos un atributo estado; por eso se proyecta antes de reunir. El mensajero se une con EMPLEADO mediante una reunión theta, porque los atributos tienen nombres distintos. La selección después de agrupar corresponde a HAVING; la de antes, a WHERE.


\subsubsection{C7. Envíos que requieren gestión}
\label{q:C7}
\textit{Usuarios:} U3. \textit{Operaciones:} selección, reunión natural, proyección y unión. Reúne en una sola lista los envíos que tienen una novedad abierta o una reclamación abierta.

\begin{center}\adjustbox{max width=\linewidth}{$\displaystyle\begin{array}{l}N \leftarrow \Pi_{\text{numero\_guia}}(\sigma_{\text{estado} = \text{'ABIERTA'}}(\text{NOVEDAD}) \bowtie \text{ASIGNACION\_DESPACHO}) \\[4pt] Q \leftarrow \Pi_{\text{numero\_guia}}(\sigma_{\text{estado} = \text{'ABIERTA'} \wedge \neg\,\text{isnull}(\text{numero\_guia})}(\text{RECLAMACION})) \\[4pt] N \cup Q\end{array}$}\end{center}

\begin{lstlisting}[style=sqlpro]
SELECT a.numero_guia
FROM   novedad n
JOIN   asignacion_despacho a ON a.id_asignacion = n.id_asignacion
WHERE  n.estado = 'ABIERTA'
UNION
SELECT numero_guia
FROM   reclamacion
WHERE  estado = 'ABIERTA'
AND    numero_guia IS NOT NULL
\end{lstlisting}
La unión elimina los repetidos: un envío con novedad y reclamación abiertas aparece una sola vez.


\subsubsection{C8. Clientes que usan todas las modalidades}
\label{q:C8}
\textit{Usuarios:} U1. \textit{Operaciones:} proyección, reunión natural y división. Encuentra los clientes que han enviado al menos un paquete en cada una de las modalidades del servicio.

\begin{center}\adjustbox{max width=\linewidth}{$\displaystyle\begin{array}{l}S \leftarrow \Pi_{\text{id\_sede, id\_modalidad}}(\text{ENVIO} \bowtie \text{SOLICITUD\_RECOLECCION}) \\[4pt] U \leftarrow \Pi_{\text{id\_cliente, id\_modalidad}}(S \bowtie \text{SEDE\_CLIENTE}) \\[4pt] D \leftarrow U \div \Pi_{\text{id\_modalidad}}(\text{MODALIDAD}) \\[4pt] \Pi_{\text{razon\_social}}(D \bowtie \text{CLIENTE})\end{array}$}\end{center}

\begin{lstlisting}[style=sqlpro]
SELECT c.razon_social
FROM   cliente c
JOIN   sede_cliente s          ON s.id_cliente   = c.id_cliente
JOIN   solicitud_recoleccion r ON r.id_sede      = s.id_sede
JOIN   envio e                 ON e.id_solicitud = r.id_solicitud
GROUP  BY c.razon_social
HAVING COUNT(DISTINCT e.id_modalidad) = (SELECT COUNT(*)
                                         FROM   modalidad)
\end{lstlisting}
ENVIO y SEDE\_CLIENTE tienen ambos el atributo id\_zona (la zona de destino y la de la sede); por eso se proyecta S antes de reunirla con SEDE\_CLIENTE. SQL no tiene un operador de división; el resultado equivalente se obtiene comparando cuántas modalidades distintas usó cada cliente con el total de modalidades.


\subsubsection{C9. Vehículos disponibles para un mensajero}
\label{q:C9}
\textit{Usuarios:} U2. \textit{Operaciones:} selección, reunión natural y proyección. Al armar un despacho, lista los vehículos disponibles del tipo que conduce el mensajero (regla RN11).

\begin{center}\adjustbox{max width=\linewidth}{$\displaystyle\begin{array}{l}M \leftarrow \sigma_{\text{id\_empleado} = \text{7}}(\text{MENSAJERO}) \\[4pt] V \leftarrow \sigma_{\text{estado} = \text{'DISPONIBLE'}}(\text{VEHICULO}) \\[4pt] \Pi_{\text{id\_vehiculo, placa, marca, modelo}}(M \bowtie V)\end{array}$}\end{center}

\begin{lstlisting}[style=sqlpro]
SELECT v.id_vehiculo, v.placa, v.marca, v.modelo
FROM   mensajero m
JOIN   vehiculo v ON v.id_tipo_vehiculo = m.id_tipo_vehiculo
WHERE  m.id_empleado = 7
AND    v.estado = 'DISPONIBLE'
\end{lstlisting}
La reunión natural une M y V por su único atributo común, id\_tipo\_vehiculo.


\subsubsection{C10. Estado actual de cada envío}
\label{q:C10}
\textit{Usuarios:} U1, U2, U3. \textit{Operaciones:} agregación, renombre, reunión natural, selección y proyección. Obtiene el estado actual de cada envío, que es el de su evento más reciente.

\begin{center}\adjustbox{max width=\linewidth}{$\displaystyle\begin{array}{l}U \leftarrow \rho_{\text{U(numero\_guia, ultima)}}\big({}_{\text{numero\_guia}}\mathcal{G}_{\text{max}(\text{fecha\_hora})}(\text{EVENTO\_ENVIO})\big) \\[4pt] R \leftarrow \sigma_{\text{fecha\_hora} = \text{ultima}}(\text{EVENTO\_ENVIO} \bowtie U) \\[4pt] \Pi_{\text{numero\_guia, nombre}}(R \bowtie \text{ESTADO\_ENVIO})\end{array}$}\end{center}

\begin{lstlisting}[style=sqlpro]
SELECT ev.numero_guia, es.nombre AS estado_actual
FROM   evento_envio ev
JOIN   estado_envio es ON es.id_estado = ev.id_estado
WHERE  ev.fecha_hora = (SELECT MAX(x.fecha_hora)
                        FROM   evento_envio x
                        WHERE  x.numero_guia = ev.numero_guia)
ORDER  BY ev.numero_guia
\end{lstlisting}
El resultado intermedio U corresponde en SQL a una subconsulta correlacionada. Es la misma lógica de la vista v\_estado\_actual\_envio.


\subsection{Tablas, Relaciones, Consultas, Informes}

Esta sección especifica lo que contendrá la base de datos del módulo y
cómo la usará cada grupo de usuarios: sus tablas, las relaciones entre
ellas, las consultas y transacciones que se ejecutarán y los informes
que se entregarán. Su construcción y ejecución corresponden al
prototipo funcional (sección~\ref{sec:prototipo}).

\subsubsection{Tablas}

La base de datos del módulo tiene las 19 tablas del modelo relacional
(figura~\ref{fig:dmrelacional}), con las columnas del diccionario de
datos (tabla~\ref{tab:diccionario}) y los objetos físicos del
capítulo 5. La tabla~\ref{tab:tablas6} las clasifica según su papel e
indica quién registra sus datos. Los catálogos y las tablas de
referencia los carga el administrador de la base de datos al
instalarla, y los actualiza a solicitud de las áreas dueñas mientras se
integran los módulos M2 a M4 (tabla~\ref{tab:interfaces}).

{\singlespacing\small
\begin{longtable}{L{3.6cm} L{1.8cm} L{5.0cm} L{4.0cm}}
\caption{Tablas de la base de datos del módulo.}
\label{tab:tablas6}\\
\toprule
\textbf{Tabla} & \textbf{Tipo} & \textbf{Contenido} & \textbf{Quién registra los datos} \\
\midrule
\endfirsthead
\toprule
\textbf{Tabla} & \textbf{Tipo} & \textbf{Contenido} & \textbf{Quién registra los datos} \\
\midrule
\endhead
\bottomrule
\endfoot
centro\_distribucion & Maestra & Centro donde se reciben, clasifican y despachan los envíos. & Centro (U2) \\
zona & Maestra & Sector de la ciudad atendido por un centro. & Centro (U2) \\
cliente & Referencia & Empresa que contrata el servicio. & Administrador \\
sede\_cliente & Referencia & Punto de recolección de un cliente. & Administrador \\
modalidad & Catálogo & Tipo de servicio según el tiempo de entrega. & Administrador \\
tipo\_vehiculo & Catálogo & Catálogo de tipos de vehículo. & Administrador \\
vehiculo & Referencia & Vehículo de la flota. & Administrador \\
empleado & Referencia & Persona vinculada a la empresa. & Administrador \\
mensajero & Referencia & Empleado que sale a ruta (subtipo de EMPLEADO). & Administrador; el centro (U2) cambia su zona \\
destinatario & Maestra & Persona que recibe los envíos. & Mesa de servicio (U3), al registrar el envío \\
despacho & Operación & Salida de un mensajero con un vehículo. & Centro (U2) \\
solicitud\_recoleccion & Operación & Pedido de recolección en una sede del cliente. & Mesa de servicio (U3); el mensajero (U4) la marca atendida \\
envio & Operación & Paquete o documento identificado por su guía. & Mesa de servicio (U3); el centro (U2) agrega peso real y flete \\
asignacion\_despacho & Operación & Cada intento de un envío en un despacho. & Centro (U2); el mensajero (U4) registra el resultado \\
estado\_envio & Catálogo & Catálogo de estados del envío. & Administrador \\
evento\_envio & Operación & Cada cambio de estado de un envío. & Centro (U2), mesa (U3) y mensajero (U4) \\
prueba\_entrega & Operación & Constancia de recibo de una entrega o devolución. & Mensajero (U4) \\
novedad & Operación & Situación que impidió la entrega en un intento. & Mensajero (U4); la mesa (U3) la gestiona \\
reclamacion & Operación & Inconformidad de un cliente. & Mesa de servicio (U3) \\
\end{longtable}
\par}

\textit{Catálogo}: valores fijos que usan las reglas y las vistas.
\textit{Referencia}: datos maestros que mantendrán los módulos M2 a M4.
\textit{Maestra}: datos propios del módulo que cambian poco.
\textit{Operación}: registros que se crean a diario con el trabajo de la
empresa.

\subsubsection{Relaciones}

Las relaciones del modelo son las 27 claves foráneas de la
tabla~\ref{tab:conversion}. En la base de datos se implementan con
restricciones de clave foránea, que el gestor verifica en cada
operación:

\begin{itemize}[itemsep=2pt]
  \item \textbf{Al insertar o modificar} una fila hija, la fila padre a
        la que apunta debe existir. Las claves foráneas son obligatorias,
        salvo tres que admiten nulos porque la relación es opcional: la
        zona del mensajero, el despacho de recolección de una solicitud
        todavía no atendida y la guía de una reclamación que no trata
        sobre un envío.
  \item \textbf{Al intentar borrar} una fila padre que tiene filas
        hijas, el gestor rechaza la operación (acción restrictiva). No se
        usa el borrado en cascada, porque ningún registro se elimina: se
        anula o se cierra por estado (RN16).
  \item \textbf{Las claves primarias no se modifican}: las sustitutas
        las genera el gestor y el número de guía no cambia una vez
        impreso, así que no se necesita actualización en cascada.
\end{itemize}

La tabla~\ref{tab:relaciones6} muestra, para cada tabla, a qué tablas
apunta y qué tablas apuntan a ella; es la guía para saber qué tablas
reunir en una consulta.

{\singlespacing\footnotesize
\begin{longtable}{L{3.5cm} L{5.4cm} L{5.5cm}}
\caption{Relaciones de cada tabla (entre paréntesis, la clave foránea).}
\label{tab:relaciones6}\\
\toprule
\textbf{Tabla} & \textbf{Apunta a (tablas padre)} & \textbf{Es referenciada por (tablas hijas)} \\
\midrule
\endfirsthead
\toprule
\textbf{Tabla} & \textbf{Apunta a (tablas padre)} & \textbf{Es referenciada por (tablas hijas)} \\
\midrule
\endhead
\bottomrule
\endfoot
centro\_distribucion & -- & zona (id\_centro); vehiculo (id\_centro); empleado (id\_centro); despacho (id\_centro) \\
zona & centro\_distribucion (id\_centro) & sede\_cliente (id\_zona); mensajero (id\_zona); envio (id\_zona) \\
cliente & -- & sede\_cliente (id\_cliente); reclamacion (id\_cliente) \\
sede\_cliente & cliente (id\_cliente); zona (id\_zona) & solicitud\_recoleccion (id\_sede) \\
modalidad & -- & envio (id\_modalidad) \\
tipo\_vehiculo & -- & vehiculo (id\_tipo\_vehiculo); mensajero (id\_tipo\_vehiculo) \\
vehiculo & centro\_distribucion (id\_centro); tipo\_vehiculo (id\_tipo\_vehiculo) & despacho (id\_vehiculo) \\
empleado & centro\_distribucion (id\_centro) & mensajero (id\_empleado); evento\_envio (id\_empleado) \\
mensajero & empleado (id\_empleado); zona (id\_zona); tipo\_vehiculo (id\_tipo\_vehiculo) & despacho (id\_mensajero) \\
destinatario & -- & envio (id\_destinatario) \\
despacho & centro\_distribucion (id\_centro); vehiculo (id\_vehiculo); mensajero (id\_mensajero) & solicitud\_recoleccion (id\_despacho); asignacion\_despacho (id\_despacho) \\
solicitud\_recoleccion & sede\_cliente (id\_sede); despacho (id\_despacho) & envio (id\_solicitud) \\
envio & solicitud\_recoleccion (id\_solicitud); destinatario (id\_destinatario); zona (id\_zona); modalidad (id\_modalidad) & asignacion\_despacho (numero\_guia); evento\_envio (numero\_guia); reclamacion (numero\_guia) \\
asignacion\_despacho & envio (numero\_guia); despacho (id\_despacho) & prueba\_entrega (id\_asignacion); novedad (id\_asignacion) \\
estado\_envio & -- & evento\_envio (id\_estado) \\
evento\_envio & envio (numero\_guia); estado\_envio (id\_estado); empleado (id\_empleado) & -- \\
prueba\_entrega & asignacion\_despacho (id\_asignacion) & -- \\
novedad & asignacion\_despacho (id\_asignacion) & -- \\
reclamacion & cliente (id\_cliente); envio (numero\_guia) & -- \\
\end{longtable}
\par}

\subsubsection{Consultas}

La tabla~\ref{tab:consultas6} resume las consultas de la
sección~\ref{sec:algebra}: quién las usa, qué datos indica y con qué
frecuencia. A ellas se suman las consultas de las vistas de usuario
(tabla~\ref{tab:vistas}), que alimentan los informes.

{\singlespacing\footnotesize
\begin{longtable}{L{0.7cm} L{2.9cm} L{1.2cm} L{1.9cm} L{4.4cm} L{2.2cm}}
\caption{Catálogo de consultas del módulo.}
\label{tab:consultas6}\\
\toprule
\textbf{Id} & \textbf{Consulta} & \textbf{Usu.} & \textbf{Datos de entrada} & \textbf{Tablas} & \textbf{Frecuencia} \\
\midrule
\endfirsthead
\toprule
\textbf{Id} & \textbf{Consulta} & \textbf{Usu.} & \textbf{Datos de entrada} & \textbf{Tablas} & \textbf{Frecuencia} \\
\midrule
\endhead
\bottomrule
\endfoot
C1 & Datos de un envío & U3 & Número de guía & ENVIO & Cada llamada de un cliente \\
C2 & Recorrido de un envío & U3, U9 & Número de guía & EVENTO\_ENVIO,\allowbreak\ ESTADO\_ENVIO,\allowbreak\ EMPLEADO & Cada consulta de seguimiento \\
C3 & Planilla de un despacho & U2, U4 & Despacho & ASIGNACION\_DESPACHO,\allowbreak\ ENVIO,\allowbreak\ DESTINATARIO,\allowbreak\ ZONA & Al salir y al cerrar cada despacho \\
C4 & Envíos clasificados sin asignar & U2 & Ninguna & EVENTO\_ENVIO,\allowbreak\ ESTADO\_ENVIO,\allowbreak\ ASIGNACION\_DESPACHO & Al armar las rutas \\
C5 & Envíos por zona y modalidad & U1 & Fecha inicial y final & ENVIO,\allowbreak\ ZONA,\allowbreak\ MODALIDAD & Semanal y mensual \\
C6 & Mensajeros con varias novedades & U1 & Fecha inicial y final & NOVEDAD,\allowbreak\ ASIGNACION\_DESPACHO,\allowbreak\ DESPACHO,\allowbreak\ EMPLEADO & Mensual \\
C7 & Envíos que requieren gestión & U3 & Ninguna & NOVEDAD,\allowbreak\ ASIGNACION\_DESPACHO,\allowbreak\ RECLAMACION & Diaria \\
C8 & Clientes que usan todas las modalidades & U1 & Ninguna & ENVIO,\allowbreak\ SOLICITUD\_RECOLECCION,\allowbreak\ SEDE\_CLIENTE,\allowbreak\ MODALIDAD,\allowbreak\ CLIENTE & Trimestral \\
C9 & Vehículos disponibles para un mensajero & U2 & Mensajero & MENSAJERO,\allowbreak\ VEHICULO & Al armar cada despacho \\
C10 & Estado actual de cada envío & U1, U2, U3 & Ninguna & EVENTO\_ENVIO,\allowbreak\ ESTADO\_ENVIO & Permanente (tableros) \\
\end{longtable}
\par}

Además de consultar, los usuarios registran la operación mediante las
transacciones de la tabla~\ref{tab:transacciones}. Cada una agrupa
varias inserciones y actualizaciones que se ejecutan como una unidad:
si alguna falla, por ejemplo porque un disparador detecta que se
incumple una regla de negocio, se deshacen todas, y la base de datos
queda como estaba (atomicidad y consistencia de las transacciones).

{\singlespacing\small
\begin{longtable}{L{3.6cm} L{1.0cm} L{7.2cm} L{2.6cm}}
\caption{Transacciones del módulo.}
\label{tab:transacciones}\\
\toprule
\textbf{Transacción} & \textbf{Usu.} & \textbf{Operaciones sobre las tablas} & \textbf{Reglas verificadas} \\
\midrule
\endfirsthead
\toprule
\textbf{Transacción} & \textbf{Usu.} & \textbf{Operaciones sobre las tablas} & \textbf{Reglas verificadas} \\
\midrule
\endhead
\bottomrule
\endfoot
Registrar solicitud de recolección & U3 & Inserta SOLICITUD\_RECOLECCION. & RN15 \\
Registrar envío & U3 & Inserta DESTINATARIO si no existe, ENVIO y su EVENTO\_ENVIO REGISTRADO (procedimiento de registro). & RN01 a RN04, RN12 \\
Recolectar & U4 & Actualiza la SOLICITUD\_RECOLECCION a ATENDIDA e inserta el evento RECOLECTADO de cada envío. & RN12 \\
Recibir y clasificar & U2 & Inserta los eventos EN\_CENTRO y CLASIFICADO y actualiza en ENVIO el peso real y el flete. & RN03, RN12 \\
Armar el despacho de entrega & U2 & Inserta DESPACHO, una ASIGNACION\_DESPACHO por envío y el evento EN\_RUTA de cada uno. & RN07, RN09 a RN12 \\
Registrar una entrega & U4 & Inserta PRUEBA\_ENTREGA, actualiza el resultado de la asignación e inserta el evento ENTREGADO. & RN05, RN06, RN08, RN12 \\
Registrar una novedad & U4 & Inserta NOVEDAD, actualiza el resultado de la asignación e inserta el evento CON\_NOVEDAD. & RN07, RN08, RN12 \\
Registrar una devolución & U4 & Igual que una entrega, en el tercer intento y con resultado y evento DEVUELTO. & RN05, RN07, RN12 \\
Cerrar el despacho & U2 & Actualiza en DESPACHO la hora de cierre y el estado CERRADO. & -- \\
Gestionar una novedad & U3 & Actualiza en NOVEDAD la acción, la solución y el estado. & -- \\
Registrar y responder una reclamación & U3 & Inserta RECLAMACION y luego actualiza su respuesta, fecha de respuesta y estado. & RN14 \\
Anular un envío & U3 & Inserta el evento ANULADO (solo antes de salir a ruta). & RN12, RN16 \\
\end{longtable}
\par}

\subsubsection{Informes}

Los informes son las salidas que cada usuario necesita para su trabajo
(sección~\ref{sec:eps}). Cada uno se obtiene de una vista de usuario
(tabla~\ref{tab:vistas}), que ya contiene las reuniones entre tablas,
filtrada y ordenada según el caso. La tabla~\ref{tab:informes} define
el contenido de cada informe.

{\singlespacing\footnotesize
\begin{longtable}{L{2.2cm} L{1.0cm} L{2.9cm} L{3.4cm} L{2.0cm} L{1.8cm}}
\caption{Informes del módulo.}
\label{tab:informes}\\
\toprule
\textbf{Informe} & \textbf{Usu.} & \textbf{Vista de origen} & \textbf{Contenido} & \textbf{Filtro y orden} & \textbf{Frecuencia} \\
\midrule
\endfirsthead
\toprule
\textbf{Informe} & \textbf{Usu.} & \textbf{Vista de origen} & \textbf{Contenido} & \textbf{Filtro y orden} & \textbf{Frecuencia} \\
\midrule
\endhead
\bottomrule
\endfoot
Planilla de ruta & U4, U2 & v\_planilla\_despacho & Orden de visita, guía, destinatario, teléfono, dirección, zona, frágil, alto valor y resultado. & Un despacho; por orden de visita. & Cada despacho \\
Seguimiento del envío & U9, U3 & v\_seguimiento\_envio & Fecha y hora, estado, observación y responsable de cada evento. & Una guía (U9: solo las de su empresa); por fecha. & Cada consulta \\
Envíos por estado & U1, U2 & v\_estado\_actual\_envio & Estado actual y número de envíos. & Por centro o zona; agrupado por estado. & Diaria \\
Cumplimiento de entregas & U1 & v\_cumplimiento\_entregas & Centro, zona, entregados, entregados a tiempo y porcentaje a tiempo. & Un centro; de menor a mayor porcentaje. & Semanal y mensual \\
Envíos por despachar & U2 & v\_envios\_por\_despachar & Guía, zona, modalidad, fecha de compromiso, peso, frágil y alto valor. & Un centro; por fecha de compromiso. & Antes de cada salida \\
Novedades abiertas & U3 & v\_novedades\_abiertas & Guía, intento, tipo y descripción de la novedad, destinatario, cliente y teléfonos. & Todas; por hora de la visita. & Diaria \\
Reclamaciones pendientes & U1, U3 & v\_reclamaciones\_pendientes & Fecha de radicación, cliente, guía, tipo y descripción. & Todas; de la más antigua a la más reciente. & Diaria \\
\end{longtable}
\par}

\subsection{Prototipo funcional de la Base de Datos}
\label{sec:prototipo}

\subsubsection{Arquitectura del prototipo}
\label{sec:arquitectura}

El prototipo se organiza como una arquitectura en cuatro capas lógicas.
Cada capa tiene una sola responsabilidad y solo usa la capa que está
debajo de ella:

\begin{enumerate}[itemsep=2pt]
  \item \textbf{Presentación.} Las aplicaciones de la
        tabla~\ref{tab:aplicaciones}: muestran formularios, planillas e
        informes, capturan los datos y validan solo su formato. No
        deciden reglas del negocio.
  \item \textbf{Casos de uso.} Un servicio por cada grupo de casos de
        uso del módulo (recolección, centro, ruta, servicio al cliente,
        indicadores y administración). Cada operación verifica el rol
        del usuario, abre la transacción, pide al dominio que aplique las
        reglas, manda guardar y confirma o deshace todo
        (tabla~\ref{tab:transacciones}).
  \item \textbf{Modelo de dominio.} Las entidades del modelo conceptual
        convertidas en clases con comportamiento, agrupadas alrededor
        del envío (con sus eventos, destinatario, prueba y novedades) y
        del despacho (con sus asignaciones). Aquí se aplican las reglas
        de negocio de la tabla~\ref{tab:reglasnegocio}, como las
        transiciones de estado (tabla~\ref{tab:transiciones}) o el
        máximo de intentos, antes de tocar la base de datos.
  \item \textbf{Datos.} Los repositorios que leen y guardan las entidades
        mediante JDBC, con las consultas de la sección~\ref{sec:algebra}
        y las vistas de usuario, y la base de datos Oracle con las
        restricciones, disparadores, índices y roles del capítulo 5.
\end{enumerate}

La presentación se ejecuta en los equipos de los usuarios; las capas de
casos de uso, dominio y acceso a datos, en el servidor de aplicaciones;
y la base de datos, en su propio servidor (sección~\ref{sec:almacenamiento}).
Las reglas de negocio se verifican dos veces: en el dominio, para
explicar al usuario qué falló antes de intentar guardar, y en la base de
datos, como garantía final aunque alguien llegue a ella por otro camino.

\textbf{Por qué no Oracle REST Data Services.} Oracle ofrece una
herramienta, Oracle REST Data Services (ORDS), que publica directamente
las tablas, vistas y procedimientos de la base de datos como servicios
web, sin escribir un servidor de aplicaciones. Se descartó como base del
prototipo porque las aplicaciones hablarían directamente con la base de
datos y se saltarían las capas de casos de uso y de dominio, lo que trae
estos problemas:

\begin{itemize}[itemsep=2pt]
  \item \textbf{Transacciones incompletas.} Operaciones como registrar
        una entrega insertan la prueba, actualizan la asignación e
        insertan el evento. Sin una capa que las agrupe, la aplicación
        tendría que hacer varias llamadas independientes; si la conexión
        del mensajero se cae en medio, la operación queda a medias.
  \item \textbf{Reglas repartidas.} Para dar mensajes claros, cada
        aplicación (web, móvil y portal) tendría que repetir las reglas
        de negocio por su cuenta, con el riesgo de que difieran entre sí
        y de la base de datos.
  \item \textbf{Aplicaciones atadas a las tablas.} Los servicios
        reflejarían la estructura de las tablas, así que cualquier cambio
        del esquema, como las columnas que se agregarán al integrar el
        módulo de facturación, obligaría a cambiar todas las
        aplicaciones.
  \item \textbf{Control de acceso más pobre.} Los permisos quedarían a
        nivel de tabla y no de caso de uso: es difícil expresar, por
        ejemplo, que un mensajero solo puede registrar resultados en sus
        propios despachos, y quedaría expuesta a internet una superficie
        mayor de la base de datos.
\end{itemize}

ORDS podría usarse más adelante para publicar consultas de solo lectura,
como el seguimiento del portal, sin afectar esta arquitectura.

El resto del prototipo se desarrollará en la siguiente entrega: se
creará la base de datos del módulo en el gestor, con los objetos físicos
del capítulo 5; se cargarán datos de prueba; y se ejecutarán las
consultas de la sección~\ref{sec:algebra}, las transacciones de la
tabla~\ref{tab:transacciones} y los informes de la
tabla~\ref{tab:informes}, para mostrar sus resultados.
